\chapter{Future Enhancement and Left Off}

\section{Current Status}

The OpenROAD integration with eSim was successfully completed and validated using the \textbf{FullAdder} and \textbf{HalfAdder} example circuits. The implemented workflow successfully performs automatic Verilog generation, OpenROAD Flow Scripts execution, backend file generation, and layout visualization using both OpenROAD GUI and KLayout.

The RTL-to-GDSII flow was tested end-to-end and the generated backend outputs were verified successfully during implementation.

\section{Completed Work}

At the end of this internship, the following implementation tasks were completed successfully:

\begin{itemize}
    \item OpenROAD execution button added to the eSim graphical interface.
    
    \item \texttt{netlist2rtl.py} implemented for automatic Verilog generation.
    
    \item \texttt{OpenROAD.py} backend module developed for ORFS execution.
    
    \item Automatic generation of \texttt{config.mk} and \texttt{constraint.sdc}.
    
    \item Successful RTL-to-GDSII flow execution for FullAdder and HalfAdder.
    
    \item Automatic OpenROAD GUI launch after flow completion.
    
    \item Successful GDSII visualization using KLayout.
    
    \item Multiple backend integration and execution issues resolved.
\end{itemize}

\section{Support for Additional Circuit Types}

The current implementation primarily supports combinational circuits such as FullAdder and HalfAdder. Future work can extend the parser to support:

\begin{itemize}
    \item sequential circuits,
    \item flip-flop based designs,
    \item hierarchical netlists,
    \item larger digital modules.
\end{itemize}

Additional improvements may also be required for handling more complex SPICE netlists generated from larger projects.

\section{Improved RTL Parsing}

The current \texttt{netlist2rtl.py} module uses a simplified parsing approach for extracting ports and subcircuit instances from the generated \texttt{.cir.out} files.

Future improvements may include:

\begin{itemize}
    \item improved port identification,
    \item better wire extraction,
    \item support for complex net connections,
    \item automatic detection of input and output signals.
\end{itemize}

A more robust parser would improve compatibility with larger and more complex circuit designs.

\section{Improved GUI Integration}

The current OpenROAD flow executes as a backend subprocess launched directly from eSim.

Future improvements may include:

\begin{itemize}
    \item real-time progress display,
    \item flow execution logs inside the GUI,
    \item flow cancellation support,
    \item automatic error reporting,
    \item generated file summaries after execution.
\end{itemize}

These features would improve the overall usability of the integration.

\section{Automatic File Management}

Currently, generated backend outputs such as \texttt{.def}, \texttt{.gds}, and \texttt{.spef} files are stored inside the ORFS results directory.

Future improvements may include:

\begin{itemize}
    \item automatic copying of generated files into the eSim project directory,
    \item automatic workspace refresh,
    \item organized output management for multiple projects.
\end{itemize}

This would simplify project management and improve integration consistency.

\section{Enhanced Layout Visualization}

The generated layouts are currently visualized using external OpenROAD GUI and KLayout applications.

Future improvements may include direct integration of layout visualization inside the eSim graphical interface, allowing users to inspect generated layouts without leaving the eSim environment.

Such integration would provide a more unified open-source RTL-to-GDSII workflow within a single application.